\begin{proposition}[E.30 --- Proved: all three remaining rational coordinates lie in the ambient-metabelian kernel]
\label{E-remaining-kernel}
Retain the actual covering inclusion $N\hookrightarrow P$ of E.26,
with compatible tangential base points, and put
\[
 C=N/\gamma_4N,\qquad
 \bar C=N/(P''\gamma_4N),\qquad
 A=\gamma_3N/\gamma_4N.
\]
Let $W\subset V=N^{\rm ab}$ be the primitive rank-eight
cusp lattice. The kernel on the third lower central graded piece is
the actual Galois-stable submodule
\[
 K_3=\ker\bigl(A\longrightarrow\operatorname{gr}_3\bar C\bigr)
       =[W,\operatorname{Lie}_2(V)].
 \tag{E.30.1}
\]
It is saturated of rank $304$ in $A$, which has rank $330$.
The rational coefficient module carrying the three coordinates of
E.29 lies entirely in $K_{3,\mathbb Q_3}$. Consequently
\[
 [\eta]\longmapsto0\quad\text{in}\quad
 H^1\bigl(G_{\mathbb Q,S_3},(A/K_3)_{\mathbb Q_3}(-1)\bigr),
 \tag{E.30.2}
\]
and $[\eta]$ has a unique preimage in
$H^1(G_{\mathbb Q,S_3},K_{3,\mathbb Q_3}(-1))$.
Thus the specified once-punctured metabelian comparison annihilates
the entire rational class retained in E.29. This assertion evaluates
its image under that comparison. Propositions~\ref{E-ross-edge}
and~\ref{E-kummer-values} below identify its elliptic coordinate
with an exact cup-product edge and evaluate its two Kummer coordinates.
\end{proposition}
\begin{proof}
We first verify the exact kernel calculation in the supplied
LM7 tower return, E.50, rather than treating its table as an input.
Choose two handle generators $e_1,e_2$ and eight cusp generators
$t_1,\ldots,t_8$ for $N$. Such a basis is obtained by cutting the
nine-sheeted torus cover along its two handle loops and its cusp
paths; the remaining ninth cusp is eliminated by the one surface
relation. The actual filling map $N\to T_3E$ kills the $t_i$.
The composite $N\to P^{\rm ab}$ is three times that map, since the
proper covering is $[3]:E\to E$. Multiplication by three is
injective on $T_3E$, so
\[
 P'=\ker(N\longrightarrow P^{\rm ab})
    =\overline{\langle\!\langle t_1,\ldots,t_8,[e_1,e_2]
                              \rangle\!\rangle}_N.
\]
This is also a direct presentation: quotienting by those generators
leaves the free abelian pro-three group on $e_1,e_2$, mapping
isomorphically to $3P^{\rm ab}$.

In $C$, write $J$ for the image of $P'$, $T$ for the subgroup
generated by the eight cusp generators, and $C_2=\gamma_2C$.
Then $J=T C_2$, the group $C_2$ is abelian, and $[T,C_2]$
is central. Commutator collection gives
\[
 J'=[T,T]\,[T,C_2].
\]
The degree-two images of $[t_i,t_j]$ are independent basis
vectors of $\bigwedge^2W\subset\bigwedge^2V$.
It follows that $[T,T]\cap\gamma_3C=\gamma_3T$;
this latter subgroup is contained in $[T,C_2]$.
Therefore
\[
 J'\cap\gamma_3C=[T,C_2]=[W,\operatorname{Lie}_2(V)].
\]
Since $J'$ is the image of $P''$, this proves (E.30.1)
as an equality of actual kernels, including at coefficient prime three.

Here is an integral basis check. Number the two handles $0,1$
and the cusp generators $2,\ldots,9$. Use the Hall basis
$[[v_i,v_j],v_k]$ with $i>j$ and $k\ge j$.
The subgroup in (E.30.1) has as its basis exactly the $304$
Hall vectors whose last index $k$ is a cusp index.
Each is directly the negative of $[v_k,[v_i,v_j]]$.
Conversely every such bracket reduces to these vectors:
when $k<j<i$, Jacobi gives
\[
 [[v_i,v_j],v_k]
   =[[v_i,v_k],v_j]-[[v_j,v_k],v_i],
\]
and both new last indices are cusp indices whenever $k$ is.
The $26$ retained vectors are the two pure handle triples
and the $24$ vectors
$[[t,e_1],e_1]$, $[[t,e_1],e_2]$, $[[t,e_2],e_2]$.
This proves saturation and the ranks over $\mathbb Z_3$.
The independent program \texttt{remaining\_kernel.py} checks
all $360$ generating brackets and all $304$ individual basis
witnesses over the integers; it uses no program absent from the
supplied return.

Pass now to the actual rational compact splitting
$V_{\mathbb Q_3}=H\oplus W_{\mathbb Q_3}$.
The kernel consists of all degree-three terms containing at
least two boundary letters, together with
$[W_{\mathbb Q_3},\bigwedge^2H]$.
Indeed Jacobi expresses every bracket with two boundary letters
in $[W,\operatorname{Lie}_2V]$, and the remaining one-boundary
generators are exactly the latter displayed subspace.

E.29 places the class in the image of cohomology of the
Galois-stable direct sum
\[
 B^{\mu_6}_{\mathbb Q_3}
       \ \oplus\
 (A_{\mathbb Q_3}^{E[3]\rtimes\mu_6})_{HW^2}
       \ \subset\ A_{\mathbb Q_3},
\]
with twist $(-1)$ understood. The first summand is
$[W^{\mu_6}_{\mathbb Q_3},w]$, where
$w\in\bigwedge^2H$; the second has two boundary letters.
Both lie in the kernel just computed. This proves (E.30.2).
The quotient $A/K_3$ has rational summands
$\operatorname{Lie}_3(H)$ and
$\operatorname{Sym}^2H\otimes W$, of Frobenius weights
three and four before the twist. After twist $(-1)$ their
weights are one and two, so a good Frobenius has no eigenvalue
one. Thus $H^0(G_{\mathbb Q,S_3},(A/K_3)_{\mathbb Q_3}(-1))=0$.
The long exact cohomology sequence makes the map from
$H^1(K_{3,\mathbb Q_3}(-1))$ injective and identifies its image
with the kernel in (E.30.2). This proves the asserted uniqueness.

This result specifies the additional kernel data that a
theta-function or elliptic path construction must retain to
evaluate E.29. It proves no vanishing before the quotient map.
The anchors are \S51.16, folios 733v--736v, scans 172--178,
and task LM7-E.23(i).
\end{proof}
